Mullin's enumeration of tree-rooted maps \cite{Mullin} gives the
combinatorial starting point: the selected spanning tree and its
complementary dual tree encode a map by a two-coordinate walk. The
continuum counterpart is the mating-of-trees construction of
Duplantier, Miller, and Sheffield \cite{DMS}. At $\gamma=\sqrt2$ its
boundary-length coordinates are independent Brownian motions. Miller
and Sheffield \cite{MSSphere} give the corresponding finite-area
sphere description, and their imaginary-geometry construction
\cite{IG4} supplies the space-filling $\mathrm{SLE}_8$ between the two
continuum trees. Couplings of discrete and continuum encodings have
also been used to compare graph distances across random-map models
\cite{GHS}. The continuum metric itself is characterized in
\cite{GMMetric}, with conformal covariance established in
\cite{GMConformal}.

Liouville Brownian motion was constructed by Garban, Rhodes, and Vargas
\cite{GRV} and independently by Berestycki \cite{Berestycki} through a
change of Brownian time determined by quantum area. Its Dirichlet-form
description \cite{GRVForm} separates two structures: the energy is
conformally invariant, while the speed measure is quantum area. This
distinction explains why the clock is a separate part of a random-map
scaling limit.

Several works establish important parts of the discrete picture.
Gwynne, Miller, and Sheffield prove convergence of the Tutte embedding
of mated-CRT maps, including conditional walk-law convergence modulo
time parameterization \cite{GMS}. Their estimates for harmonic
functions compare discrete and continuum energies up to constants and control
oscillation \cite{GMSHarmonic}; their invariance principle for ergodic
scale-free environments develops a broader framework for embedded
networks \cite{GMS2}. Gwynne and Miller obtain resistance,
return-probability, and displacement estimates for several infinite
random planar maps, including spanning-tree-weighted maps \cite{GM}.

Berestycki and Gwynne prove parameterized, quenched convergence to
Liouville Brownian motion for mated-CRT maps \cite{BG}. Their
deterministic time scale is an annealed exit-time median, comparable
to the inverse quantum area of a cell. Whether its ratio to that
inverse area converges is left open in \cite[Remark~1.3]{BG}.
Theorem~\ref{thm:main} concerns
the finite combinatorial tree-weighted map and its stationary
half-edge walk. It retains the conditional path law together with
metric and area, and identifies the acceleration as exactly $n$
for \eqref{eq:continuum-form}.

Our geometric inputs are the companion results in \cite{MetricMaps}.
The nondegeneration argument adapts the local-flow method in
\cite[Section~``The canonical conformal structure'']{ConformalMaps}
to a primal and a dual electrical network. The use of planar duality
to determine an effective conductivity has classical precedents,
including Dykhne's exact-conductivity argument \cite{Dykhne}. Here
we must first construct the energy limit without regularity
assumptions on the shapes of discrete faces. A separate Green
estimate in the original finite-map law then controls sources of
small stationary mass and supplies the missing clock information.
